% BHU PYQ -- Manually transcribed & error-corrected
% Paper: GLB-301: Petrology and Economic Geology
% Exam: B.Sc. (Hons.) Semester III Examination, 2022-23

\documentclass[11pt,a4paper]{article}
\usepackage[a4paper,top=2.5cm,bottom=2.5cm,left=2.8cm,right=2.8cm,headheight=14pt]{geometry}
\usepackage{amsmath,amssymb}
\usepackage[T1]{fontenc}
\usepackage[utf8]{inputenc}
\usepackage{lmodern,microtype,fancyhdr,enumitem,hyperref,lastpage,parskip}

\hypersetup{colorlinks=true,urlcolor=black,linkcolor=black,
  pdftitle={GLB-301 Petrology and Economic Geology -- BHU B.Sc. Sem III 2022-23}}
\pagestyle{fancy}\fancyhf{}
\renewcommand{\headrulewidth}{0.4pt}\renewcommand{\footrulewidth}{0.4pt}
\fancyhead[L]{\small\textit{Banaras Hindu University}}
\fancyhead[R]{\small\textit{GLB-301: Petrology \& Economic Geology}}
\fancyfoot[C]{\small Page \thepage\ of \pageref{LastPage}}

\newlist{parts}{enumerate}{2}
\setlist[parts,1]{label=(\alph*),leftmargin=*,itemsep=5pt,topsep=3pt}
\newcommand{\pts}[1]{\hfill\textit{[#1]}}

\begin{document}
\begin{center}
  {\large\bfseries BANARAS HINDU UNIVERSITY}\\[2pt]
  {\normalsize B.Sc. (Hons.) --- Semester III Examination, 2022-23}\\[6pt]
  \rule{0.8\linewidth}{0.5pt}\\[6pt]
  {\large\bfseries Paper: GLB-301 --- Petrology and Economic Geology}\\[4pt]
  {\normalsize Time: 3 Hours \hfill Maximum Marks: 70}\\[2pt]
  {\small Ref. No.: 055019}
\end{center}
\rule{\linewidth}{0.4pt}
\smallskip
\noindent\textit{Instructions: Answer \textbf{five} questions in all, selecting at least \textbf{two} each from Sections A and B, and Section-C which is compulsory.}
\bigskip

\bigskip\noindent{\large\bfseries SECTION --- A}
\rule{\linewidth}{0.4pt}\smallskip

\noindent\textbf{Question 1.} Write in brief about the classification of igneous rocks. \pts{14}

\medskip\noindent\textbf{Question 2.} Discuss about primary sedimentary structures with suitable sketch. \pts{14}

\medskip\noindent\textbf{Question 3.} Briefly describe the types of metamorphism. Discuss the principal controlling agents of different types of metamorphism. \pts{14}

\medskip\noindent\textbf{Question 4.} Write short notes on \textbf{any four} of the following: \pts{$4 \times 3\frac{1}{2} = 14$}
\begin{parts}
  \item Porphyritic texture
  \item Dyke and Sill
  \item Limestone
  \item Sandstone
  \item Metamorphic zones
  \item Isogrades
\end{parts}

\bigskip\noindent{\large\bfseries SECTION --- B}
\rule{\linewidth}{0.4pt}\smallskip

\noindent\textbf{Question 5.} Write important chemical formula, distribution in India and uses of \textbf{any four} of the following economic minerals: \pts{$4 \times 3\frac{1}{2} = 14$}
\begin{parts}
  \item Baryte
  \item Corundum
  \item Chalcopyrite
  \item Gypsum
  \item Talc
  \item Coal
\end{parts}

\medskip\noindent\textbf{Question 6.} Write in brief \textbf{any two} of the following: \pts{$2 \times 7 = 14$}
\begin{parts}
  \item What is coal? Discuss about the origin and economic uses of coal.
  \item What are refractory minerals? Discuss about the economic uses of refractory minerals and their distribution in India.
  \item What is chromite? Describe about economic uses and distribution of chromite ore deposits in India.
\end{parts}

\medskip\noindent\textbf{Question 7.} Briefly describe about classification of economic mineral deposits. \pts{14}

\bigskip\noindent{\large\bfseries SECTION --- C}
\rule{\linewidth}{0.4pt}\smallskip

\noindent\textbf{Question 8.} Answer the following:
\begin{parts}
  \item Write the names of minerals formed in discontinuous Bowen's reaction series. \pts{2}
  \item Differentiate between breccia and conglomerate. \pts{2}
  \item Define cementation of sedimentary rocks. \pts{1}
  \item What is protolith? \pts{1}
  \item Write the name of a high grade metamorphic rock. \pts{1}
  \item Write the names of two petroleum deposits in India. \pts{2}
  \item Write the name of one ore mineral of aluminium. \pts{1}
  \item How many sets of cleavage planes are there in fluorite? \pts{1}
  \item Which metal is extracted from pyrolusite? \pts{1}
  \item What are the economic uses of galena? \pts{2}
\end{parts}

\vfill\rule{\linewidth}{0.4pt}
\begin{center}\small
\href{https://bhu-exam.vercel.app/}{Click here to visit our website}\\[4pt]
\href{https://me-aryan.vercel.app/}{Click here to visit our contributor}\\[4pt]
\href{https://quashan-me.vercel.app/}{Click here to visit our contributor}
\end{center}
\end{document}
